\part{Differential Geometry}
\chapter{Smooth Manifolds}
\begin{importantdefinition}{Smooth Manifolds via Charts}{smoothmanifoldcharts}
	An \tib{$n$-dimensional smooth manifold} is a second-countable Hausdorff space $M$ together with a collection of maps called "charts" such that:
	\begin{enumerate}
		\item Each chart is a homeomorphism $\Phi : U \to U' \subseteq \bR^n$, where $U$ is open in $M$ and $U'$ is open in $\bR^n$.
		\item Each point $x \in M$ is in the domain of some chart,
		\item Given charts $\Phi : U \to U'$ and $\Psi : V \to V'$, the coordinate change function $\Phi\Psi^{-1} : \Psi(U \cap V) \to \Phi(U \cap V)$ is smooth (i.e. infinitely continuously differentiable)
		\item The collection of charts is maximal, i.e. if a map fulfills conditions (1) and doesn't violate condition (3) when taken together with preexisting charts, then that new map is a chart.
	\end{enumerate}
	We call our any collection of charts satisfying conditions $1,2$ and $3$ an \tib{atlas} of $M$.
\end{importantdefinition}
Note that this definition directly generalizes topological manifolds, but with additional smoothness demands on our charts. Since most manifolds can't be equipped with the structure of a vector space, differentiability isn't well-defined for our charts $\Phi : U \to U \subseteq \bR^n$, so instead we move that demand to the coordinate change functions $\Phi\Psi^{-1}$, which are functions mapping $\bR^n$ into itself.
\newpar
We also want to construct a second, slightly more abstract definition of a smooth manifold, based on the notion of \tib{functional structures}.
\begin{importantdefinition}{}{}
	Let $(X, \cO)$ be a topological space. A \tib{fuctional structure} on $X$ is a function $F_X$ defined on the collection $\cO$ of open sets in $X$, such that:
	\begin{enumerate}
		\item $F_X(U)$ is a subalgebra of the algebra $C(U, \bR)$ of continuous real-valued functions on $U$$^1$\marginnote{$^1$Originally defined in Exercise \ref{ex:realfunctionalgebra}.},
		\item $F_X(U)$ contains all constant functions,
		\item If $V \subseteq U$ and $f \in F_X(U)$, then $f|_V \in F_X(V)$, and
		\item If $U = \bigcup_{i \in I} U_i$ and $f|_{U_i} \in F_x(U_i)$ for all $i$, then $f \in F_X(U)$.
	\end{enumerate}
	The pair $(X, F)$ is called a \tib{functionally structured space}.
\end{importantdefinition}
\noindent
In particular, we care about the following functional structures on $X = \bR^n$:
\begin{enumerate}
	\item $F_X(U) = C(U, \bR)$,
	\item $F_X(U) = C^k(U, \bR)$,
	\item $F_X(U) = C^\infty(U, \bR)$,
	\item $F_X(U) = A(U, \bR)$ ($=$ the set of analytic functions $U \to \bR$).
\end{enumerate}
\begin{exercise}
	If $(X, F_X)$ is a functionally structured space, we can turn $(U, F_U)$ into a functionally structured space by setting $F_U(V) = F_X(V)$ for all subsets $V \subseteq U$.
\end{exercise}
\begin{importantdefinition}{}{}
	A \tib{morphism of functionally structured spaces}
	\begin{align*}
		(X, F_X) \to (Y, F_Y)
	\end{align*}
	is a map $\Phi : X \to Y$ such that precomposition pulls back elements of $F_Y(U)$ to elements of $F_X(\Phi^{-1}(U))$:
	\begin{align*}
		f \in F_Y(U) \implies f \circ \Phi \in F_X(\Phi^{-1}(U))
	\end{align*}
	An isomorphism, as always, is a morphism whose inverse exists and is a morphism.
\end{importantdefinition}
\begin{importantdefinition}{Smooth Manifolds via Functional Structure}{smoothmanifoldfunctional}
	An $n$-dimensional differentiable manifold is a second countable \ul{functionally structured} Hausdorff space $(M, F)$ which is locally isomorphic to $(\bR^n, C^\infty)$, i.e. each point in $M$ has a neighborhood $U$ such that $(U, F_U) \cong (V, C^\infty_V)$.
\end{importantdefinition}
\pagebreak
\begin{theorem}
	Our two definitions of smooth manifolds are equivalent.
\end{theorem}
\begin{proof}
	\begin{enumerate}
		\item[$\Rightarrow:$] Let $M$ be a smooth manifold in the sense of Definition \ref{df:smoothmanifoldcharts}. We need to define the sets $F(U)$ for all open sets $U$. Let $U$ be the domain of a chart $\Phi : U \to U' \subseteq \bR^n$. Then we can set
		\begin{align*}
			F(U) = \lr{f \circ \Phi \mid f \in C^\infty(U', \bR)}
		\end{align*}
		Then $F$ inherits the properties defining a functional structure from $C^\infty(U', \bR)$, and $\Phi$ is an isomorphism between $F$ and $C^\infty$ by definition of $F$.
		\item[$\Leftarrow:$] Let $(M, F)$ be a smooth manifold in the sense of Definition \ref{df:smoothmanifoldfunctional}. We take our charts $\Phi : U \to V$ to be isomorphisms $(U, F_U) \cong (V, C_V^\infty)$. 
		\newpar
		It remains to be shown that the coordinate change function $\theta := \Phi\Psi^{-1}$ of any pair of charts is $C^\infty$, and that our atlas is maximal.
		Since $\theta$ is a morphism, $f \circ \theta$ is $C^\infty$ for any $f \in C^\infty$. Therefore, $\pi_i \circ \theta$ is $C^\infty$ for all $i \in [1..n]$, so $\theta$ is also $C^\infty$. 
		\newpar
		At the same time, if $\Phi$ and $\Psi$ are homeomorphisms and $\theta$ is $C^\infty$, then $\theta$ is clearly an isomorphism, since $\theta \in C^\infty$ implies $f \circ \theta \in C^\infty$ for all $f \in C^\infty$ and $\theta^{-1} \in C^\infty$ exists and fulfills the same conditions. Therefore, our atlas is maximal.
	\end{enumerate}
\end{proof}
\begin{proposition}
	These two constructions are inverse to each other.
\end{proposition}
From now on, we will implicitly assume that every functional structure induced by a chart structure and every chart structure induced by a functional structure is created via this construction, where charts correspond to local isomorphisms of the functional structures.
\newpar
This lets us generalize smoothness from vector spaces to manifolds:
\begin{importantdefinition}{}{}
	A map between manifolds is called \tib{smooth} or \tib{$C^\infty$} if it is a morphism between their respective functional structures
\end{importantdefinition}
\begin{importanttheorem}{}{}
	A map $f : M \to N$ between two smooth manifolds is smooth in the sense of being a morphism between functional structures if and only if, for any charts $\Phi : U \to U' \subseteq \bR^m$ on $M$ and $\Psi : V \to V' \subseteq \bR^n$ on $N$, the function $\Psi \circ f \circ \Phi^{-1}$ is smooth.
\end{importanttheorem}
\begin{proof}
	\begin{enumerate}
		\item[$\Rightarrow$:] Let $f : (M, F_M) \to (N, F_N)$ be a morphism. Let $\pi_i : \bR^n \to \bR$ be the $i$-th coordinate projection function. Since charts are isomorphisms, we get:
		\begin{align*}
			&\pi_i \in C^\infty(\bR^n, \bR)\\
			\implies &\pi_i \circ \Psi \in F_N(V)\\
			\implies &\pi_i \circ \Psi \circ f \in F_M(f^{-1}(V))\\
			\implies &\pi_i \circ \Psi \circ f \circ \Phi^{-1} \in C^\infty(\Phi(f^{-1}(V)), \bR)
		\end{align*}
		Since every component of $\Psi \circ f \circ \Phi^{-1}$ is smooth, the entire function is smooth.
		
		
		\item[$\Leftarrow$:] Let $\Psi \circ f \circ \Phi^{-1}$ be smooth for all charts $\Phi, \Psi$. Since $\Phi$ is bijective, the set on which this function is well-defined is 
		\begin{align*}
			\Phi(f^{-1}(V)) \cap \Phi(U) = \Phi(f^{-1}(V) \cap U).
		\end{align*}
		Let $g \in F_N(U)$. Then we have $g = h \circ \Psi$ for some chart $\Psi : V \to \bR^n$ and some smooth real-valued function $h \in C^\infty(\bR^n, \bR)$.
		\newpar
		We get:
		\begin{align*}
			&g := h \circ \Phi \in F_N(V)\\
			\implies &h \in C^\infty(\bR^n, \bR)\\
			\implies &h \circ \Psi \circ f \circ \Phi^{-1} \in C^\infty(\Phi(f^{-1}(V) \cap U),\bR)\\
			\implies &(g \circ f)|_{f^{-1}(V) \cap U} = h \circ \Psi \circ f \circ \Phi^{-1} \circ \Phi \in F_M(f^{-1}(V) \cap U, \bR)
		\end{align*}
		Picking different charts $\Phi_i : U_i \to \bR^n$ whose domains cover $f^{-1}(V)$ gives $g \circ f \in F_M(f^{-1}(V))$ as desired.
	\end{enumerate}
\end{proof}
Another concept that manifolds inherit from vector spaces is orientation:
\begin{importantdefinition}{}{}
	Let $M$ be an $n$-manifold. Then an atlas $A$ which is maximal with regards to the property that, for any pair of charts $\Phi, \Psi$ in the atlas, the Jacobian of the coordinate change function $\Phi\Psi^{-1}$ has positive determinant at all points in its domain, is called an \tib{orientation} of $M$. An $n$-manifold having such an atlas is called \tib{orientable}.
\end{importantdefinition}
\begin{theorem}
	Every orientable manifold has at least two orientations.
\end{theorem}
\begin{proofsketch}
	Given an orientation of $M$, we can pick one coordinate of each chart and reverse the sign, leading to a new orientation.
\end{proofsketch}
\begin{theorem}
	Every connected manifold has at most two orientations.
\end{theorem}
\begin{proof}
	Since charts are bijective, their Jacobians never have determinant zero.
	\newpar
	Now, $\cA$ and $\cB$ be two orientations. Define a map
	$s(p) : M \to \lr{1,-1}$ by setting $s(p) = 1$ if all coordinate change functions $\Phi\Psi^{-1}$ involving one chart from each orientation have positive determinant (in which case all of them have positive determinant, since $\Phi\psi^{-1} = \Phi\Psi^{-1}\Psi\psi^{-1}$), and $s(p) = -1$ otherwise. Since Jacobian determinants vary continuously, $s(p)$ is locally constant. Since $M$ is connected, $s(p)$ must be constant, leaving two options:
	\begin{enumerate}
		\item If $s(p) = 1$, then $\cA$ and $\cB$ are compatible, and must thus already contain each others charts, since orientations are maximal.
		\item If $s(p) = -1$, let $\cC$ be another orientation for which the analogously defined function $s_C$ comparing $\cC$ to $\cA$ is negative. Then, given charts $\Phi_A, \Phi_B, \Phi_C$ from each orientation whose domains match, we get
		\begin{align*}
			\Phi_B \Phi_C^{-1} = (\Phi_B \Phi_A^{-1}) \circ (\Phi_A \Phi_C^{-1})
		\end{align*}
		since the Jacobians of $(\Phi_B \Phi_A^{-1})$ and $(\Phi_A \Phi_C^{-1})$ have negative signs, the Jacobian of their composition must have positive sign, so $\cB$ and $\cC$ must be compatible, and thus identical.
	\end{enumerate}
\end{proof}
\begin{definition}
	As with topological manifolds, smooth manifolds with boundary are defined by taking smooth manifolds and replacing the domain of our charts with the upper half-space $\bH^n := \lr{x \in \bR^n \mid x_n \geq 0}$.
\end{definition}